\section{Discussion and Limits}\label{sec:discussion}
\paragraph{What a designer can take from this.} First, a fixed ranking is a bet on the fault rate. The per-bit term of Eq.~\eqref{eq:p-expansion} rewards rankings that keep the sign of a flip's effect when faults are sparse; the pairwise term rewards balanced, sign-blind rankings when they are dense; and the rate at which one takes over can be computed before any fault is injected. \Branch{R1}{On ResNet-8 that computation located the switch within a factor of 1.5 of where it was measured on every checkpoint, about 12\% late on average.}{\TBD{How far the computed switch can be trusted, given Stage 1.}} Second, when the rate is not known, the absolute gradient is the safe static choice: it never fell far behind in the static evaluation, while the positive gradient lost up to 18 points where faults were dense and significance-first protection 9 to 34 points at a moderate rate. Third, \Branch{R2}{when the rate is known or bounded, the rate-aware assignment removes the bet: on ResNet-8 it was never more than half a point behind the best tested mask \TBD{}.}{\TBD{what the rate-aware assignment does and does not buy, from Stage 2}.} Fourth, protection should be judged by its residual error curve, not by its name. Codes and checksum correction keep a quadratic residual only while a second error per word or channel is unlikely; beyond that they protect almost nothing, while copies keep their power law (Fig.~\ref{fig:residual}). At dense fault rates, triplicating ranked bits beat SEC-DED codes and checksum correction of equal storage by 30 points or more on ResNet-8.

\paragraph{Why not search for the mask by fault injection?} A Monte Carlo search estimates the same expected loss by sampling and needs many fault fields to resolve differences of a fraction of a point; on the pinned ResNet-8 such a search and a relaxed per-bit optimizer did not beat the absolute-gradient mask where the model is usable, and the relaxed optimizers lost 4 to 12 points there (Appendix~\ref{app:search}). The expression of Section~\ref{sec:theory} computes the expectation in closed form, to second order around the average faulty weights, and samples only the part it approximates. \Branch{R2c}{The sampled correction helped where \TBD{}.}{The sampled correction did not help on ResNet-8.}

\paragraph{Limits.} The theory is a second-order description around the average faulty weights and predicts loss, not accuracy. The confirmatory tests cover nine ResNet-8 checkpoints, and one further model, the DS-CNN, is checked with a smaller design (\TBD{Stage 4 outcome}); the other four models enter through the static evaluation only. The prediction places the switch about 12\% late on ResNet-8, and below 6.25\% protection the switch falls where the model is no longer usable, so the prediction is tested only from 6.25\% up. Each model except ResNet-8 has one trained checkpoint, several input cohorts had been used before, and the derivation was written after the reversal of Appendix~\ref{app:calibration} had been observed, although every prediction tested in Section~\ref{sec:switch} was fixed before its fault fields were drawn. The residual curves assume independent faults, fault-free voters and decoders, and no miscorrection of three or more errors; the checksum decoder is applied to stored weights, which favours it, and an ABFT that checks computed outputs is not modelled. Counting storage bits ignores the area, power and latency of voters, decoders and their placement. Biases, scales, embeddings and normalization parameters are never faulted, and activation, datapath and control faults are outside the fault model. Logical $p$ is a model parameter, not a radiation rate, and the band in which masks differ, about $10^{-4}$ to $10^{-1}$, is a regime of dense memory errors that scrubbing would shorten~\cite{biswas2005avf}.
